\documentclass[10pt]{article}

\usepackage[margin=1in]{geometry}
\usepackage{amsmath}
\usepackage{amssymb}
\usepackage{amsthm}
\usepackage{mathtools}
\usepackage[shortlabels]{enumitem}
\usepackage[normalem]{ulem}
\usepackage{booktabs}
\usepackage{tcolorbox}

\usepackage{hyperref}
\hypersetup{
  colorlinks   = true, %Colours links instead of ugly boxes
  urlcolor     = black, %Colour for external hyperlinks
  linkcolor    = blue, %Colour of internal links
  citecolor    = blue  %Colour of citations
}

\usepackage{courier}
\usepackage{listings}
\lstset{numbers=left}
\lstset{basicstyle=\ttfamily}

%%%%%%%%%%%%%%%%%%%%%%%%%%%%%%%%%%%%%%%%%%%%%%%%%%%%%%%%%%%%%%%%%%%%%%%%%%%%%%%%

\theoremstyle{definition}
\newtheorem{problem}{Problem}
\newcommand{\E}{\mathbb E}
\newcommand{\R}{\mathbb R}
\DeclareMathOperator{\Var}{Var}
\DeclareMathOperator*{\argmin}{arg\,min}
\DeclareMathOperator*{\argmax}{arg\,max}

\newcommand{\trans}[1]{{#1}^{T}}
\newcommand{\loss}{\ell}
\newcommand{\w}{\mathbf w}
\newcommand{\mle}[1]{\hat{#1}_{\textit{mle}}}
\newcommand{\map}[1]{\hat{#1}_{\textit{map}}}
\newcommand{\normal}{\mathcal{N}}
\newcommand{\x}{\mathbf x}
\newcommand{\y}{\mathbf y}
\newcommand{\ltwo}[1]{\lVert {#1} \rVert}

\newcounter{ProblemCounter}
\newcommand{\nextproblem}{\filbreak\stepcounter{ProblemCounter} \noindent\textbf{Problem \theProblemCounter}:}

%%%%%%%%%%%%%%%%%%%%%%%%%%%%%%%%%%%%%%%%%%%%%%%%%%%%%%%%%%%%%%%%%%%%%%%%%%%%%%%%

\begin{document}

\begin{center}
{
\Large
    Quiz (Transactions and Locks II)
}

    \vspace{0.1in}
\end{center}

\noindent
\textbf{Instructions:}
For each problem, you must perform all of the following tasks correctly in order to get credit.
\begin{enumerate}
\item Circle all SQL commands that block.
\item Write an X through all commands that error.
\item If a command causes a deadlock, write \lstinline{DEADLOCK} and stop the problem.
\item Write the output of each \lstinline{SELECT} statement.
\end{enumerate}
\textbf{Rules:}
\begin{enumerate}
    \item You are allowed to use a computer and access the internet.
        (I strongly encourage you to refer to the postgres documentation and/or class notes.)
    \item You may not use an LLM.
    \item You may not connect to a running postgres database.
\end{enumerate}

\vspace{1in}

\nextproblem
\begin{tcolorbox}
\begin{tabular}{p{3.25in}p{3in}}
\hspace{-0.2in}Session 1
\begin{lstlisting}
CREATE TABLE t (a INT PRIMARY KEY);
CREATE TABLE u (b INT UNIQUE);
BEGIN;
INSERT INTO t VALUES (1);
INSERT INTO u VALUES (1);
LOCK TABLE t IN SHARE MODE;
SELECT count(*) FROM u;
INSERT INTO u VALUES (2);

COMMIT;
SELECT count(*) FROM t;
\end{lstlisting}
    &
\hspace{-0.2in}Session 2
\begin{lstlisting}





BEGIN ISOLATION LEVEL
REPEATABLE READ;
SELECT count(*) FROM u;
INSERT INTO t VALUES (2);

INSERT INTO u VALUES (2);
COMMIT;
SELECT count(*) FROM u;
\end{lstlisting}
\end{tabular}
\vspace{-0.2in}
\end{tcolorbox}
\vspace{0.25in}

%%%%%%%%%%%%%%%%%%%%%%%%%%%%%%%%%%%%%%%%
\nextproblem
\begin{tcolorbox}
\begin{tabular}{p{3.25in}p{3in}}
\hspace{-0.2in}Session 1
\begin{lstlisting}
CREATE TABLE t ( a INT UNIQUE );
BEGIN;
INSERT INTO t VALUES (1);
INSERT INTO t VALUES (2);
INSERT INTO t VALUES (3);


ABORT;
  
\end{lstlisting}
    &
\hspace{-0.2in}Session 2
\begin{lstlisting}




BEGIN;
INSERT INTO t VALUES (3);


INSERT INTO t VALUES (1);
COMMIT;
SELECT count(*) FROM t;
\end{lstlisting}
\end{tabular}
\vspace{-0.2in}
\end{tcolorbox}
\vspace{0.25in}

%%%%%%%%%%%%%%%%%%%%%%%%%%%%%%%%%%%%%%%%
\nextproblem
\begin{tcolorbox}
\begin{tabular}{p{3.25in}p{3in}}
\hspace{-0.2in}Session 1
\begin{lstlisting}
CREATE TABLE t(a INT UNIQUE);
CREATE TABLE u(b INT REFERENCES t(a));
INSERT INTO t VALUES (8);
INSERT INTO t VALUES (9);
BEGIN;

INSERT INTO u VALUES (9);

COMMIT;
\end{lstlisting}
    &
\hspace{-0.2in}Session 2
\begin{lstlisting}






BEGIN;
UPDATE t SET a=10 WHERE a=9;

COMMIT;
\end{lstlisting}
\end{tabular}
\vspace{-0.2in}
\end{tcolorbox}
\vspace{0.25in}

%%%%%%%%%%%%%%%%%%%%%%%%%%%%%%%%%%%%%%%%
\nextproblem
\begin{tcolorbox}
\begin{tabular}{p{3.25in}p{3in}}
\hspace{-0.2in}Session 1
\begin{lstlisting}
CREATE TABLE t (a INT UNIQUE);
CREATE TABLE u (
    b INT REFERENCES t(a)
    DEFERRABLE INITIALLY DEFERRED
    );
BEGIN;

INSERT INTO t VALUES (8);
ABORT;
\end{lstlisting}
    &
\hspace{-0.2in}Session 2
\begin{lstlisting}







BEGIN;

INSERT INTO u VALUES (8);
COMMIT;
\end{lstlisting}
\end{tabular}
\vspace{-0.2in}
\end{tcolorbox}
\vspace{0.25in}

\end{document}
